\documentclass[11pt]{article}
\usepackage[margin=0.8in]{geometry}
\usepackage[T1]{fontenc}
\usepackage{lmodern,amsmath,amssymb,amsthm,booktabs,array}
\usepackage[colorlinks=true,linkcolor=blue,urlcolor=blue,citecolor=blue]{hyperref}
\hypersetup{pdftitle={An orphan-count identity, a composition asymptotic, and nine power-sum patterns},pdfauthor={Blake Stone},pdfsubject={Complete proofs and Lean verification of selected sequence assertions}}
\setlength{\parindent}{0pt}
\setlength{\parskip}{0.4em}
\newtheorem{theorem}{Theorem}[section]
\newtheorem{lemma}[theorem]{Lemma}
\newtheorem{corollary}[theorem]{Corollary}
\newcommand{\N}{\mathbb N}
\newcommand{\Z}{\mathbb Z}
\begin{document}
\begin{center}
{\LARGE\bfseries An orphan-count identity, a composition\\[3pt]
asymptotic, and nine power-sum patterns}\par
\vspace{0.5em}
{\large Blake Stone}\par
Proof developed with OpenAI Astra 6.2\par
November 26, 2026
\end{center}

\textbf{Abstract.}
We prove the remaining weighted-row conjecture for the recursive triangle
A389524, connecting its original recurrence to Nathanson's literal
orphan-counting definition and completing the entry's three conjectures.
We formally verify the A377824 minimum-part-position asymptotic with
explicit corrections and an exponentially small relative error. A concurrent
repository manuscript established that result before publication of this
package and is credited. We also classify nine power-sum families in an
A289629 comment: eight are true and one is false. Four true families were
already established in A349989. The additional assertion count is therefore
six: five proofs and one refutation, alongside the composition formalization.
Four new Lean modules and two reproduced dependencies verify the results
from the original definitions. Classical ingredients are explicitly credited;
bounded prior-work searches do not establish exhaustive historical priority.

\begin{center}
\begin{tabular}{@{}ll>{\raggedright\arraybackslash}p{3.7in}@{}}
\toprule
Entry & Contribution & Result\\
\midrule
A389524 & 1 proved & Weighted rows count the actual orphans; all three entry conjectures are now proved.\\
A377824 & 1 formalized & Exact main term and geometric error bound; a concurrent repository proof is credited.\\
A289629 & 4 proved, 1 refuted & Complete nine-branch classification, with four other true branches credited as prior work.\\
\bottomrule
\end{tabular}
\end{center}

\section{A389524: the weighted row counts orphans}

Mats Granvik's A389524~\cite{A389524} defines an integer triangle by the
source program's recurrence. Its first column is $T(1,1)=1$ and
$T(n,1)=0$ for $n\ne1$. For $k\ge2$, set $T(n,k)=0$ when $n<k$, and
otherwise set
\begin{equation}\label{eq:triangle7}
T(n,k)=1+\sum_{i=1}^{k-1}T(n-i,k-1)
          -\sum_{i=1}^{n-1}T(n-i,k).
\end{equation}
The exceptional first column is retained throughout the formalization.

Let $h(n)$ be Nathanson's orphan-counting function A274628~\cite{A274628}:
the number of integer quadruples satisfying
\[
ab-cd=n,\qquad a>c\ge0,\qquad b>d\ge0.
\]
Write $d(n)$ for the number of positive divisors of $n$, and
$\sigma(n)=\sum_{r\mid n}r$ for their sum.

\begin{theorem}\label{thm:weighted7}
For every positive integer $n$,
\[
\boxed{\sum_{k=1}^{n}kT(n,k)=h(n).}
\]
Equivalently, in the source's reversed ordering,
\[
\sum_{j=1}^{n}(n-j+1)T(n,n-j+1)=h(n).
\]
Both sides also have the divisor-statistic expression
\begin{equation}\label{eq:orphan-convolution7}
2h(n)+d(n)=\sum_{t=1}^{n-1}d(t)d(n-t)+3\sigma(n).
\end{equation}
\end{theorem}

This is Conjecture 2 of the entry. The earlier package proves
$\sum_kT(n,k)=d(n)$ and $T(2n,n)=n-d(n)$, its Conjectures 1 and 3.
Those earlier results are not counted again here.

\subsection{A direct count of the orphans}

We first prove~\eqref{eq:orphan-convolution7} from the literal quadruples.
They form a finite set: writing
$ab-cd=a(b-d)+d(a-c)=b(a-c)+c(b-d)$ shows $a,b\le n$, and then
$c,d<n$.

If $c=0$, the equation is $ab=n$ and $d$ may be any of
$0,\ldots,b-1$. Summing these $b$ choices over divisors $b$ of $n$
gives $\sigma(n)$ orphans.

If $c>0$, use the bijection
\[
(a,b,c,d)\longmapsto(x,y,z,w)=(a-c,b,c,b-d).
\]
Its image consists exactly of positive quadruples with
$xy+zw=n$ and $y\ge w$. The inverse is
$(a,b,c,d)=(x+z,y,z,y-w)$. Positivity and the order condition verify
both directions.

There are
\[
C(n)=\sum_{t=1}^{n-1}d(t)d(n-t)
\]
positive ordered quadruples with $xy+zw=n$: choose $t=xy$, then a
factor pair of $t$ and a factor pair of $n-t$. Swapping the two products
interchanges $y>w$ and $y<w$. The equality case $y=w$ has
$(x+z)y=n$; for each divisor $r=x+z$ there are $r-1$ positive splits.
Its count is therefore $\sigma(n)-d(n)$. It follows that
\[
h(n)=\sigma(n)+\frac{C(n)+\sigma(n)-d(n)}2,
\]
which proves~\eqref{eq:orphan-convolution7}. This classical convolution
formula is proved here as a definition bridge, rather than assumed from
the OEIS generating function.

\subsection{The weighted triangle identity}

Define $R_0(x)=0$ and
\begin{equation}\label{eq:R7}
R_{k+1}(x)=1+(1-x^k)R_k(x).
\end{equation}
The column polynomials derived from~\eqref{eq:triangle7} are
\begin{equation}\label{eq:column7}
\sum_{n\ge0}T(n,k)x^n=x^kR_k(x),\qquad k\ge1.
\end{equation}
Indeed, the first column is $x$. Moving the last sum in the recurrence
to the left gives
\[
\frac{C_k(x)}{1-x}=\frac{x^k}{1-x}
 +(x+\cdots+x^{k-1})C_{k-1}(x),
\]
and hence $C_k=x^k+x(1-x^{k-1})C_{k-1}=x^kR_k$.
The reproduced Lean module proves this bridge by induction on the
original recurrence using finite coefficient sums.

Fix $N\ge1$ and work in the quotient ring
$A_N=\Z[x]/(x^{N+1})$, writing $q$ for the class of $x$.
Every $1-q^r$ with $r\ge1$ is a unit, with a finite geometric inverse.
Define
\[
D=\sum_{r=1}^{N}\frac{q^r}{1-q^r},\qquad
E=\sum_{r=1}^{N}\frac{r q^r}{1-q^r},\qquad
W=\sum_{k=1}^{N}kq^kR_k(q).
\]
We prove the finite identity
\begin{equation}\label{eq:weighted7}
\boxed{2W=D^2-D+3E\quad\text{in }A_N.}
\end{equation}
Classical weighted divisor identities underlie this
formula~\cite{Uchimura,Bressoud,Agarwal}.
We give all the finite algebraic details below.

First introduce a marker $z$ by $J_0(z)=0$ and
\[
J_{k+1}(z)=z\bigl(1+(1-q^k)J_k(z)\bigr).
\]
Then $J_k(1)=R_k(q)$. For $d\ge1$, induction gives
\[
[z^d]J_k(z)=
\begin{cases}
\displaystyle\prod_{j=k-d+1}^{k-1}(1-q^j),&k\ge d,\\
0,&k<d.
\end{cases}
\]
We claim
\begin{equation}\label{eq:marked7}
\sum_{k=0}^{N}q^kJ_k(z)
=\sum_{d=1}^{N}\frac{q^d}{1-q^d}z^d.
\end{equation}
To prove it, put $B_t^{(d)}=\prod_{j=0}^{d-1}(1-q^{t+j})$.
Cancellation of the common factors gives
\[
B_{t+1}^{(d)}-B_t^{(d)}
=q^t(1-q^d)B_{t+1}^{(d-1)}.
\]
The coefficient of $z^d$ on the left of~\eqref{eq:marked7} is
$\sum_{t=0}^{N}q^{t+d}B_{t+1}^{(d-1)}$; extending the sum to this
upper limit only adds zero terms. Multiplication by $1-q^d$ telescopes
to $q^d(B_{N+1}^{(d)}-B_0^{(d)})=q^d$, since these endpoint products
are respectively $1$ and $0$. Coefficients with $d>N$ vanish.
This proves~\eqref{eq:marked7}. Differentiation at $z=1$ now gives
\begin{equation}\label{eq:Jderiv7}
\sum_{k=1}^{N}q^kJ'_k(1)=E.
\end{equation}

Next define the first two finite Uchimura moments
\[
M_j=\sum_{k=0}^{N}k^j q^k\prod_{i=k+1}^{N}(1-q^i),\qquad j=1,2.
\]
They satisfy
\begin{equation}\label{eq:moments7}
M_1=D,\qquad M_2=D^2+E.
\end{equation}
For a finite proof, put
\[
H(z)=\sum_{k=0}^{N}z^kq^k\prod_{i=k+1}^{N}(1-q^i),\quad
G(z)=\prod_{i=1}^{N}(1-zq^i),\quad
P=\prod_{i=1}^{N}(1-q^i).
\]
Coefficient comparison gives
$H(qz)=(1-qz)H(z)$ and $(1-qz)G(qz)=G(z)$; the endpoint terms
vanish because $q^{N+1}=0$. Thus $K(z)=H(z)G(z)$ satisfies
$K(qz)=K(z)$. For $r>0$, its coefficient of $z^r$ is annihilated by
the unit $1-q^r$, so it is zero. The constant coefficient is $P$,
and therefore $HG=P$.

At $z=1$, $G(1)=P$ is a unit and $H(1)=1$. Logarithmic differentiation
of the finite product, or two applications of the product rule, gives
\[
\frac{G'(1)}P=-D,\qquad
\frac{G''(1)}P=D^2-\sum_{i=1}^{N}
 \left(\frac{q^i}{1-q^i}\right)^2.
\]
Differentiating $HG=P$ once and twice consequently gives
$H'(1)=D$ and
$H''(1)=D^2+\sum_i(q^i/(1-q^i))^2$.
Since $M_1=H'(1)$ and $M_2=H''(1)+H'(1)$, the remaining identity is
\[
\sum_{i=1}^{N}\frac{q^i}{(1-q^i)^2}
=\sum_{i,j=1}^{N}j q^{ij}
=\sum_{i,j=1}^{N}i q^{ij}=E.
\]
These geometric expansions are finite in $A_N$, proving~\eqref{eq:moments7}.

Finally, let $A_k=kR_k(q)-J'_k(1)$. The two defining recurrences imply
\[
A_0=0,\qquad A_{k+1}=k+(1-q^k)A_k.
\]
Induction, using the common tail-product recurrence for $M_1$ and $M_2$,
gives
\[
2A_{N+1}=N(N+1)-(M_2-M_1).
\]
Summing $A_{k+1}-A_k=k-q^kA_k$ from $k=0$ to $N$ also gives
$2\sum_{k=1}^{N}q^kA_k=N(N+1)-2A_{N+1}=M_2-M_1$.
Combining this with~\eqref{eq:Jderiv7} and~\eqref{eq:moments7} yields
\[
2W=2E+(M_2-M_1)=D^2-D+3E,
\]
as claimed.

Take $N=n$. By~\eqref{eq:column7}, $[x^n]W=\sum_{k=1}^{n}kT(n,k)$.
Factor-pair counting gives $[x^n]D=d(n)$, $[x^n]E=\sigma(n)$, and
$[x^n]D^2=C(n)$. Extracting this coefficient in~\eqref{eq:weighted7}
and comparing with~\eqref{eq:orphan-convolution7} proves
Theorem~\ref{thm:weighted7}. Replacing $k$ by $n-j+1$ gives the
source's ordering. All steps involve finite sums and polynomial rings;
no analytic convergence premise is needed.

\textbf{Formal scope and attribution.}
\texttt{Orphans.lean} proves the literal quadruple count by explicit
bijections, including finiteness and the factor-pair counts.
\texttt{Weighted.lean} proves the finite algebra, extracts coefficients,
and states the exact reversed-row conjecture. The reproduced
\texttt{Uchimura.lean} and \texttt{Triangle.lean} supply the earlier
first-moment identity and original recurrence bridge. The classical
weighted identities are credited; the new assertion resolution is
their fully checked connection to this specific recursive triangle.


\section{A377824: the sum of all minimum-part positions}

For a composition $c=(c_1,\ldots,c_\ell)$ of $n$ into positive integers,
define
\[
 s(c)=\sum_{\substack{1\le j\le\ell\\c_j=\min(c)}}j,
 \qquad a(n)=\sum_{c\models n}s(c),
\]
with $a(0)=0$. Thus every occurrence of the minimum is included, with
one-based positions. This is A377824; Kot\v{e}\v{s}ovec's April 19, 2025
conjecture is $a(n)\sim n^2 2^{n-5}$~\cite{A377824}.

\begin{theorem}\label{thm:minimum-positions7}
For every integer $n\ge2$,
\[
 0\le a(n)-\frac{(n+1)(n+6)2^n}{32}
 \le n^2\left(\frac74\right)^n.
\]
In particular,
\[
 \frac{a(n)}{n^2 2^{n-5}}
 =1+\frac7n+\frac6{n^2}+\varepsilon_n,
 \qquad 0\le\varepsilon_n\le32\left(\frac78\right)^n,
\]
and the conjectured asymptotic follows.
\end{theorem}

\begin{proof}
Starting from the composition $(1)$, use the two operations
\[
 L(c)=(1,c_1,\ldots,c_\ell),\qquad
 R(c)=(c_1+1,c_2,\ldots,c_\ell).
\]
They enumerate every composition of each subsequent integer exactly
once: delete the first part when it is $1$, and otherwise decrease it
by one, to obtain the unique predecessor. The two cases are disjoint.
Consequently, the number of compositions is $C_n=2^{n-1}$ for $n\ge1$,
and the number beginning with $1$ is $H_n=2^{n-2}$ for $n\ge2$.

For a composition $c$, let $o(c)$ count its ones and let $w(c)$ sum
their positions. Let $\alpha(c)$ be the indicator that its first part
is one. Directly from the two operations,
\[
\begin{array}{c|cc}
 & L(c)&R(c)\\ \hline
o & o(c)+1&o(c)-\alpha(c)\\
w & w(c)+o(c)+1&w(c)-\alpha(c)
\end{array}
\]
The left operation shifts every old position by one. The right operation
removes a marked position only when the first part was one, and that
position has weight one.

Write $O_n=\sum_{c\models n}o(c)$ and $W_n=\sum_{c\models n}w(c)$.
Summing the table gives, for $n\ge2$,
\[
 O_{n+1}=2O_n+2^{n-2},\qquad
 W_{n+1}=2W_n+O_n+2^{n-2}.
\]
Since the compositions of $2$ are $(1,1)$ and $(2)$, the initial values
are $O_2=2$ and $W_2=3$. Substitution in the recurrences and induction
give
\begin{equation}\label{eq:ones-main7}
 O_n=(n+2)2^{n-3},\qquad
 W_n=\frac{(n+1)(n+6)2^n}{32}\qquad(n\ge2).
\end{equation}

It remains to bound all compositions whose minimum exceeds one.
Let $E_n$ count compositions without ones, including the empty
composition at zero, so $E_0=1$ and $E_1=0$. Let $D_n$ count compositions
whose tail contains no one. For $n\ge1$, the two operations give
\[
 E_{n+1}=D_n,\qquad D_{n+1}=E_n+D_n.
\]
Indeed, a left child always contains a one, while a right child contains
no one precisely when its predecessor's tail contains none. For the
tail statistic, a left child has the whole predecessor as its tail,
and a right child retains the predecessor's tail. Hence
$E_{n+2}=E_{n+1}+E_n$ for $n\ge1$; the same equation at zero follows
from $E_2=1$. Since $(7/4)^2\ge7/4+1$, induction gives
$E_n\le(7/4)^n$ for all $n\ge0$.

If $c$ contains a one, then $s(c)=w(c)$. Otherwise $w(c)=0$ and,
because all parts are positive,
\[
 0\le s(c)\le1+\cdots+\ell\le\ell^2\le n^2.
\]
Summing over all exceptional compositions proves
\[
 0\le a(n)-W_n\le n^2E_n\le n^2(7/4)^n.
\]
Equation~\eqref{eq:ones-main7} proves the first assertion. Dividing by
$n^2 2^n/32$ gives the stated correction terms and error bound; every
term after the leading one tends to zero.
\end{proof}

Knopfmacher and Munagi~\cite{KnopfmacherMunagi} study minimum-part
counts, sums of their values, and the position of the first minimum.
Their minimum-one dominance argument and refined composition counts
are established background. Here the statistic sums all minimum
positions, and its main term and error are derived directly. The Lean
proof establishes exhaustive, duplicate-free enumeration and identifies
the recursive weight with the literal one-based index sum before proving
the estimate. No generating-function or asymptotic premise is assumed.
At the final publication check, a concurrent repository manuscript~\cite{Concurrent}
already proved this asymptotic with the same main term and the sharper error
$O(n^2\varphi^n)$, where $\varphi=(1+\sqrt5)/2$.
That manuscript is retained. The present package is an independently
developed formal companion, proving the explicit bound stated here.
It is not counted as an additional assertion resolution.

\section{A289629: classification of nine power-sum patterns}

Put $F_k(n)=(k+1)^n+(-k)^n$. Thompson's December 3, 2021 comment in
A289629~\cite{A289629} conjectures nonsquarefreeness in nine families.
Four branches are already covered by formulas in A349989~\cite{A349989}
and are credited below. The assertion concerns square divisibility;
the displayed $k$ need not be the least eligible value.

\begin{theorem}\label{thm:power-patterns7}
For every integer $m\ge1$, the eight affirmative rows of the following
table have the indicated square divisor. The fifth row is false at
$m=2$.
\begin{center}
\begin{tabular}{@{}rlll l@{}}
\toprule
Branch & Exponent $n$ & $k$ & Square divisor & Status\\
\midrule
1 & $22m-11$ & $2$ & $23^2$ & Proved\\
2 & $20m-6$ & $3$ & $5^2$ & Prior formula in A349989\\
3 & $20m-2$ & $3$ & $5^2$ & Prior formula in A349989\\
4 & $3^m$ & $7$ & $13^2$ & Proved\\
5 & $15m$ & $7$ & --- & False at $m=2$\\
6 & $20m-15$ & $14$ & $11^2$ & Proved\\
7 & $16m-8$ & $19$ & $17^2$ & Prior formula in A349989\\
8 & $42m-35$ & $23$ & $43^2$ & Proved\\
9 & $8m-4$ & $113$ & $17^2$ & Prior formula in A349989\\
\bottomrule
\end{tabular}
\end{center}
The fifth row becomes true when $m$ is restricted to odd integers:
$13^2\mid F_7(15m)$ for every odd $m\ge1$.
\end{theorem}

\begin{proof}
For odd $t$, the elementary factorization
\[
 X^t+Y^t=(X+Y)
 \bigl(X^{t-1}-X^{t-2}Y+\cdots+Y^{t-1}\bigr)
\]
gives $X+Y\mid X^t+Y^t$. Taking $X=(k+1)^d$ and $Y=(-k)^d$ yields
\begin{equation}\label{eq:odd-multiples7}
 F_k(d)\mid F_k(dt)\qquad(t\text{ odd}).
\end{equation}
The required seed factorizations are
\[
\begin{aligned}
F_2(11)&=23^2\cdot331,&F_3(2)&=5^2,\\
F_7(3)&=13^2,&F_{14}(5)&=11^2\cdot1831,\\
F_{19}(8)&=17^2\cdot147347969,&F_{23}(7)&=43^2\cdot639073,\\
F_{113}(4)&=17^2\cdot1148593.&&
\end{aligned}
\]
The eight exponents, in table order with the fifth omitted, are
\[
\begin{gathered}
11(2m-1),\quad2(10m-3),\quad2(10m-1),\quad3\cdot3^{m-1},\\
5(4m-3),\quad8(2m-1),\quad7(6m-5),\quad4(2m-1).
\end{gathered}
\]
Each multiplier of its seed exponent is positive and odd for $m\ge1$.
Equation~\eqref{eq:odd-multiples7} proves all eight square-divisibility
claims. The same argument with $d=3$ and $t=5m$ proves the corrected
fifth claim when $m$ is odd.

For the counterexample, $m=2$ gives
\[
\begin{aligned}
F_7(30)&=8^{30}+7^{30}=1260479379576072532986987473\\
 &=113\cdot3361\cdot25741\cdot12001921\cdot10742679301.
\end{aligned}
\]
All five factors are distinct primes, so this integer is squarefree.
Here is a compact primality certificate for the largest factor, whose
verification is also included in Lean. The Lucas criterion states that
$p$ is prime if $a^{p-1}\equiv1\pmod p$ and
$a^{(p-1)/q}\not\equiv1\pmod p$ for every prime $q\mid p-1$
~\cite{LucasLean}. Use successively
\[
\begin{array}{c|c|l}
p&a&p-1\\ \hline
35808931&3&2\cdot3^2\cdot5\cdot23\cdot17299\\
10742679301&2&2^2\cdot3\cdot5^2\cdot35808931.
\end{array}
\]
Both full powers have residue one. For the distinct prime divisors $q$
in increasing order, the proper-power residues are respectively
\[
\begin{gathered}
(35808930,\ 1193830,\ 21860597,\ 6523943,\ 6900949),\\
(10742679300,\ 5670453206,\ 4923777471,\ 9590149649).
\end{gathered}
\]
None equals one. The smaller primalities follow by trial division.
The complete factorization, primalities, modular powers, and coprimality
checks are kernel-verified, proving squarefreeness and refuting the
universally quantified fifth branch.
\end{proof}

All nine literal branches are therefore settled: eight true and one
false. Four true branches were already established in A349989, so the
contribution count here is four remaining family proofs and one
refutation. The odd-parameter correction is not counted again. No claim
that all even parameters in the fifth family give squarefree values is
made. The Lean definitions use the original integer expression and
prove the general divisibility bridge; the counterexample uses no
external factorization or primality oracle.

\section{Verification and contribution boundaries}

Each result has a standalone package with full source definitions, a
written proof, primary source snapshots, prior-work notes, pinned build
configuration, a fresh package-build log, and a separate audit and
compilation log. The four new proof modules are the A377824Formal and A289629
\texttt{Proof.lean} files and the weighted triangle's \texttt{Orphans.lean}
and \texttt{Weighted.lean}. That package additionally reproduces the two
earlier dependencies \texttt{Uchimura.lean} and \texttt{Triangle.lean}.
The latter has only a clarified historical-scope comment; its mathematical
code is unchanged from the earlier triangle proof.

\begin{center}
\begin{tabular}{@{}l>{\raggedright\arraybackslash}p{4.6in}@{}}
\toprule
Package & Main declarations checked\\
\midrule
A389524Weighted & \texttt{Nathanson.orphan\_formula};
\texttt{weighted\_rows\_nilpotent}, \texttt{weighted\_row}, and
\texttt{original\_conjecture} in \texttt{A389524Weighted}.\\
A377824Formal & \texttt{mem\_compositions}, \texttt{compositions\_nodup},
\texttt{a\_positions}, \texttt{closed\_form\_error},
\texttt{normalized\_limit}, \texttt{original\_conjecture}.\\
A289629 & \texttt{odd\_multiple\_dvd}, \texttt{eight\_true\_patterns},
\texttt{counterexample\_squarefree}, \texttt{pattern\_5\_false},
\texttt{pattern\_5\_corrected}.\\
\bottomrule
\end{tabular}
\end{center}

All packages pin Lean 4.19.0 and mathlib commit
\texttt{c44e0c8ee63ca166450922a373c7409c5d26b00b}.
The exact delivered sources passed fresh package builds and independent
compilations. The public theorem reports contain only the standard axioms
\texttt{propext}, \texttt{Classical.choice}, and \texttt{Quot.sound}.
No proof placeholders, custom axioms, \texttt{native\_decide} calls,
or unchecked numerical certificates occur. The finite primality
certificates in the disproof are checked by the Lean kernel.

With the pinned toolchain selected, reproduction in each package requires
\texttt{lake exe cache get} followed by \texttt{lake build}.
The manifests pin transitive dependencies; \texttt{SHA256SUMS.txt}
identifies the delivered files. The separate audit logs record the fresh
compilation locations and exact source hashes.

\textbf{Counting and prior work.}
This manuscript proves the remaining A389524 conjecture, formally verifies
the A377824 asymptotic, and resolves nine A289629 branch assertions. Four
of those nine already follow from explicit formulas in A349989. A concurrent
repository manuscript had proved the composition result by the final
publication check. Excluding those credited results leaves six additional
assertions: five proofs and one refutation. The previously proved triangle
identities, the composition formalization, its correction terms and error
estimate, and the repaired odd-parameter power-sum statement are not counted
as further conjecture resolutions.

The selected source comments remained conjectural in the preserved
exports. Exact-identifier and semantic searches, inspection of the linked
primary literature, and checks of public proof collections found no
earlier exact resolution of the six additional assertions. The detailed
records are in the package \texttt{source\_checks.md} files. These checks
are bounded; they do not certify exhaustive historical priority or OEIS
editorial acceptance. The contribution is the explicit resolution and
complete formal verification of the selected statements, with their
classical ingredients and previously established cases identified.

\begin{thebibliography}{11}
\raggedright
\bibitem{A389524}
M. Granvik, OEIS A389524, October 7, 2025, including Conjecture 2 and
the defining source recurrence.
\url{https://oeis.org/A389524}.

\bibitem{A274628}
OEIS A274628, Nathanson's orphan-counting function, introduced by
N. J. A. Sloane, July 7, 2016; literal integer-solution description
by D. Radcliffe, March 28, 2019.
\url{https://oeis.org/A274628}.

\bibitem{Uchimura}
K. Uchimura, An identity for the divisor generating function arising
from sorting theory, \emph{Journal of Combinatorial Theory, Series A}
\textbf{31} (1981), 131--135.
\url{https://doi.org/10.1016/0097-3165(81)90009-1}.

\bibitem{Bressoud}
D. M. Bressoud and M. V. Subbarao, On Uchimura's connection between
partitions and the number of divisors, \emph{Canadian Mathematical
Bulletin} \textbf{27} (1984), 143--145.
\url{https://doi.org/10.4153/CMB-1984-022-5}.

\bibitem{Agarwal}
A. Agarwal, S. C. Bhoria, P. Eyyunni, and B. Maji,
Bressoud--Subbarao type weighted partition identities for a generalized
divisor function, arXiv:2210.03457 (2022). See equations (1.5) and (1.9)
for the classical generalized weight identity and second moment.
\url{https://arxiv.org/abs/2210.03457}.

\bibitem{A377824}
J. T. Rascoe, OEIS A377824, November 8, 2024; asymptotic conjecture
by V. Kot\v{e}\v{s}ovec, April 19, 2025.
\url{https://oeis.org/A377824}.

\bibitem{KnopfmacherMunagi}
A. Knopfmacher and A. O. Munagi, Smallest Parts in Compositions,
in \emph{Advances in Combinatorics}, Springer (2013), 197--207.
\url{https://doi.org/10.1007/978-3-642-30979-3_11}.

\bibitem{A289629}
J.-S. Gerasimov, OEIS A289629, September 2, 2017; nine-family
conjectural comment by K. P. Thompson, December 3, 2021.
\url{https://oeis.org/A289629}.

\bibitem{A349989}
J. E. Schoenfield, OEIS A349989, December 7, 2021; established square
divisibility formulas covering four branches of the A289629 comment.
\url{https://oeis.org/A349989}.

\bibitem{LucasLean}
B. Bailey and mathlib contributors, Lucas primality theorem,
pinned mathlib source.
\url{https://github.com/leanprover-community/mathlib4/blob/c44e0c8ee63ca166450922a373c7409c5d26b00b/Mathlib/NumberTheory/LucasPrimality.lean}.

\bibitem{Concurrent}
B. Stone, Positions of minimum parts in integer compositions.
Concurrent repository manuscript, present at the final publication check
in commit \texttt{c92ab48c98ac}; developed with OpenAI Codex.
\url{https://github.com/blakestone-x/math/tree/c92ab48c98acba51cb42ceb73536a7bad251ec2d/A377824}.
\end{thebibliography}
\end{document}
